\subsection*{典型模式：1 生产者 + 1 消费者}

\co@inthreads.async@ 最常见的用途，是矩阵乘法里的 \textbf{1P1C}（一生产者、一消费者）划分：

\begin{lstlisting}[language=croqtile,caption={典型模式：1 生产者 1 消费者}]
parallel p1 by 2 : thread {
  inthreads.async (p1 == 0) {
  }
  inthreads.async (p1 == 1) {
  }
}
\end{lstlisting}

声明 \co@parallel p1 by 2 : thread@ 在同一个外层块中建立两个线程角色。\co@p1@ 的值选择职责，并不表示 K 方向的分块编号。

\co@p1 == 0@ 的区域承担生产者职责，负责写入共享暂存缓冲区。若写入使用另一个已空闲槽位，另一组就可以同时推进计算。

\co@p1 == 1@ 的区域承担消费者职责，读取已完成的分块并更新累加器。两个区域共享暂存存储，因此事件协议必须明确槽位何时从可写变为可读，以及何时重新可写。

共享地址空间要求明确的读写协调：没有协调的话，消费者可能在生产者还没写完缓冲区时就开始读。事件用于表达这一先后关系。
